\chapter{Fundamentação Teórica}
\label{cap:fundamentacao-teorica}

Este capítulo reúne os conceitos necessários para acompanhar a proposta e a sua avaliação. A Seção~\ref{sec:bpmn} apresenta a modelagem de processos e a notação BPMN 2.0, que é o formato de saída dos experimentos. A Seção~\ref{sec:llms} trata dos modelos de linguagem, de como geram texto e de por que o custo de uso depende do número de \textit{tokens}. A Seção~\ref{sec:dsl-compressao} apresenta as linguagens de domínio específico, as gramáticas formais, a transpilação e a compressão semântica, que são a base da BPMN-DSL e do transpilador. A Seção~\ref{sec:adaptacao} descreve o ajuste de modelos a uma tarefa, usado para especializar o modelo de pequeno porte. A Seção~\ref{sec:nl-bpmn} trata da extração de modelos de processo a partir de texto e das medidas usadas para avaliá-los, e a Seção~\ref{sec:estatistica}, dos métodos estatísticos usados para comparar as abordagens.

\section{Processos de Negócio e a Notação BPMN 2.0}
\label{sec:bpmn}\label{subsec:bpm}\label{subsec:bpmn-notacao}\label{subsec:bpmn-xml}

Um modelo de processo descreve quais atividades uma organização executa, em que ordem, sob quais condições e por quem. O Gerenciamento de Processos de Negócio (BPM, do inglês \textit{Business Process Management}) trata os processos de uma organização como objetos que podem ser descritos, analisados, executados e melhorados de forma sistemática \cite{dumas2018fundamentals}. A área reúne contribuições da administração, da engenharia de software e da pesquisa operacional, e ganha importância quando o processo deixa de ser uma prática informal e passa a ser um artefato que pode ser verificado \cite{weske2019business}.

Essa formalização é necessária porque a automação, a auditoria e a verificação de conformidade dependem de uma representação do processo que todos os envolvidos interpretem da mesma forma. A notação usada precisa atender a dois públicos ao mesmo tempo. Analistas e gestores devem conseguir ler o modelo sem formação técnica, e programas precisam de uma definição precisa para analisar ou executar o processo. A BPMN atende aos dois com uma notação gráfica de leitura relativamente simples, apoiada por uma especificação formal \cite{omg2011bpmn}.

A \textit{Business Process Model and Notation}, em sua versão 2.0, é mantida pelo \textit{Object Management Group} (OMG) e é a principal referência para a representação de processos de negócio \cite{omg2011bpmn}. A especificação define duas coisas: os símbolos gráficos usados pelos analistas e um modelo formal dos mesmos elementos, com uma forma de gravação em XML. Com isso, um mesmo modelo pode ser lido por pessoas, no desenho, e processado por ferramentas, no arquivo.

A especificação agrupa os elementos da notação em quatro categorias \cite{omg2011bpmn}. Os objetos de fluxo (\textit{flow objects}) são os nós do processo: eventos, atividades e \textit{gateways}. Os objetos de conexão (\textit{connecting objects}) ligam esses nós e compreendem os fluxos de sequência, os fluxos de mensagem e as associações. As raias (\textit{swimlanes}), divididas em \textit{pools} e \textit{lanes}, indicam os participantes, os setores ou os responsáveis por cada parte do processo. Os artefatos (\textit{artifacts}), como grupos, anotações textuais e objetos de dados, acrescentam informação complementar sem alterar o fluxo. Os parágrafos seguintes detalham os objetos de fluxo, que concentram a lógica do processo e são o que um modelo gerado precisa acertar.

Os eventos marcam o início, o fim ou um ponto intermediário do processo e podem ter um gatilho, como o recebimento de uma mensagem, um temporizador ou um erro. As atividades são o trabalho executado e se dividem em tarefas, que não são decompostas no modelo, e subprocessos, que contêm um fluxo próprio. As tarefas têm tipos que indicam como são executadas, como a tarefa de usuário, feita por uma pessoa com apoio de um sistema, a tarefa de serviço, feita por um sistema, e as tarefas de envio e de recebimento de mensagens.

Os \textit{gateways} controlam a divergência e a convergência do fluxo. O exclusivo segue exatamente um dos caminhos de saída, escolhido por uma condição. O paralelo segue todos os caminhos ao mesmo tempo e, na convergência, espera que todos cheguem. O inclusivo segue um ou mais caminhos, conforme as condições que forem verdadeiras. O baseado em eventos segue o caminho do primeiro evento que ocorrer. Um mesmo tipo de \textit{gateway} aparece, portanto, em dois papéis: abrindo os ramos de uma decisão e fechando esses ramos depois. Os fluxos de sequência ligam os nós dentro de um mesmo participante, e os fluxos de mensagem ligam participantes diferentes.

Com esses elementos, a BPMN representa os padrões mais comuns de processos, como decisões, execuções paralelas, repetições, esperas e trocas de mensagens entre participantes. A notação é extensa, e \citeonline{dumas2018fundamentals} observam que essa extensão exige aprendizado tanto de quem modela quanto de qualquer sistema que tente manipular os modelos automaticamente.

A especificação define a gravação de um modelo BPMN 2.0 em XML (\textit{Extensible Markup Language}), uma linguagem de marcação em que o conteúdo é organizado em elementos delimitados por \textit{tags}, aninhados uns nos outros, com atributos e com um único elemento raiz \cite{bray2008xml}. Em um arquivo BPMN, o elemento raiz \texttt{definitions} declara os espaços de nomes (\textit{namespaces}) usados e contém um ou mais processos. Cada nó do processo é um elemento com um identificador próprio, no atributo \texttt{id}. Cada fluxo de sequência é um elemento à parte, que aponta para a origem e para o destino pelos atributos \texttt{sourceRef} e \texttt{targetRef}. As coordenadas usadas para desenhar o modelo ficam em uma seção separada do arquivo, a de intercâmbio de diagramas (BPMNDI), que não altera o significado do processo \cite{omg2011bpmn}.

Um documento XML pode estar correto em dois níveis. Ele é bem formado quando segue as regras da própria linguagem, como ter todas as \textit{tags} fechadas na ordem certa e um único elemento raiz \cite{bray2008xml}. Ele é válido quando, além disso, obedece a um esquema, que define quais elementos podem aparecer, em que ordem, com quais atributos e com quais tipos de valor. O esquema da BPMN 2.0 é escrito em XSD (\textit{XML Schema Definition}), a linguagem de esquemas do W3C \cite{thompson2004xmlschema}, e é publicado junto com a especificação. A validação contra o esquema é feita por programa e tem resultado binário: o documento é aceito ou rejeitado.

O esquema verifica a estrutura, mas não o comportamento do processo. Ele exige, por exemplo, que os identificadores sejam únicos no documento, mas declara \texttt{sourceRef} e \texttt{targetRef} como nomes qualificados, e não como referências a identificadores, de modo que um fluxo que aponta para um elemento inexistente não é rejeitado. Também não verifica se todos os nós são alcançáveis a partir do início, se uma decisão tem mais de uma saída ou se um evento de fim encerra de fato o fluxo. Um arquivo aceito pelo esquema pode, assim, representar um processo malformado ou diferente do pretendido.

A gravação em XML favorece a troca de modelos entre ferramentas, mas produz arquivos muito maiores que a descrição do processo. Mesmo um processo pequeno exige declarar os espaços de nomes, um identificador para cada nó, um elemento para cada ligação e as referências de cada ligação aos nós que ela une. O mesmo identificador aparece em vários pontos do arquivo e precisa ser escrito da mesma forma em todos eles. Para uma ferramenta de modelagem, que gera o arquivo por programa, essa repetição não é um problema. Para um modelo de linguagem, que gera o arquivo símbolo a símbolo, ela significa uma saída longa, repetitiva e sensível a desvios pequenos: uma \textit{tag} esquecida torna o documento malformado, e um identificador trocado liga elementos errados.

\section{Modelos de Linguagem}
\label{sec:llms}\label{subsec:transformer}\label{subsec:tokenizacao}\label{subsec:saidas-estruturadas}

Os modelos de linguagem de grande escala (LLMs, do inglês \textit{Large Language Models}) usados nos experimentos, tanto os de grande quanto os de pequeno porte, compartilham a mesma arquitetura e o mesmo modo de gerar texto. Esta seção descreve essa arquitetura, a divisão do texto em \textit{tokens}, que é a unidade em que os provedores cobram o uso dos modelos, e as dificuldades de gerar textos que precisam obedecer a regras formais.

A arquitetura Transformer, proposta por \citeonline{vaswani2017attention}, substituiu o processamento sequencial das redes recorrentes por um mecanismo de atenção que relaciona diretamente todas as posições de uma sequência. Cada posição é representada por um vetor, do qual se obtêm uma consulta ($Q$), uma chave ($K$) e um valor ($V$). A saída da atenção é
\begin{equation}
	\text{Atenção}(Q, K, V) = \text{softmax}\left(\frac{QK^{\top}}{\sqrt{d_k}}\right) V,
\end{equation}
em que $d_k$ é a dimensão das chaves. Cada posição recebe, assim, uma combinação dos valores de todas as outras, ponderada pela semelhança entre a sua consulta e as chaves delas. Como essas operações são multiplicações de matrizes, uma sequência inteira pode ser processada em paralelo, o que permitiu treinar modelos em escala muito maior que a das arquiteturas anteriores.

Os modelos de linguagem atuais, como o GPT-3 \cite{brown2020language} e o Qwen2.5-Coder \cite{hui2024qwen25coder}, usam a parte decodificadora dessa arquitetura de forma autorregressiva. Dada uma entrada $x$, a probabilidade de uma saída $y = (y_1, \ldots, y_T)$ é decomposta em
\begin{equation}
	p(y \mid x) = \prod_{t=1}^{T} p(y_t \mid y_{<t}, x),
\end{equation}
isto é, o modelo estima cada \textit{token} a partir da entrada e dos \textit{tokens} já gerados. Na geração, o \textit{token} escolhido a cada passo é acrescentado à sequência e passa a fazer parte do contexto do passo seguinte. A escolha depende da estratégia de decodificação: na decodificação gulosa, toma-se sempre o \textit{token} mais provável; na amostragem, sorteia-se um \textit{token} segundo a distribuição. A temperatura $\tau$ controla o quanto essa distribuição é concentrada. O modelo produz, para cada \textit{token} $i$ do vocabulário, uma pontuação $z_i$, e a probabilidade de escolha é
\begin{equation}
	p_i = \frac{\exp(z_i / \tau)}{\sum_j \exp(z_j / \tau)}.
\end{equation}
Com $\tau = 1$, usa-se a distribuição do próprio modelo. Com $\tau < 1$, as diferenças entre as pontuações aumentam e a escolha se concentra nos \textit{tokens} mais prováveis, e com $\tau > 1$ a distribuição se achata e \textit{tokens} menos prováveis passam a ser sorteados com mais frequência \cite{holtzman2020curious}. No limite em que a temperatura tende a zero, toda a probabilidade fica no \textit{token} de maior pontuação, e a amostragem equivale à decodificação gulosa, mas, em serviços de terceiros, detalhes de execução ainda podem fazer o mesmo pedido produzir saídas diferentes.

Treinados sobre grandes coleções de texto, esses modelos passam a executar tarefas descritas apenas por instruções e exemplos colocados na entrada, sem alteração dos pesos, o que \citeonline{brown2020language} chamam de aprendizado em contexto. É essa capacidade que permite pedir a um modelo de grande porte que escreva em uma notação que ele nunca viu, desde que a descrição da notação esteja no \textit{prompt}. O modo de geração tem também uma consequência para formatos estruturados. O modelo gera da esquerda para a direita e não revisa o que já emitiu, de modo que um formato que exija declarar, em um ponto do texto, uma informação que só se determina adiante obriga o modelo a antecipar uma decisão que ainda não tomou.

Os modelos de linguagem não operam sobre palavras, mas sobre \textit{tokens}, que são fragmentos de texto definidos por um tokenizador. O método mais comum, o \textit{Byte-Pair Encoding} (BPE), começa com caracteres isolados e une repetidamente os pares mais frequentes no texto de treino, até atingir o tamanho de vocabulário desejado \cite{sennrich2016neural}. Palavras comuns da língua viram um único \textit{token}, e sequências raras são quebradas em vários. Identificadores gerados por programa, nomes de atributos e símbolos de marcação tendem a cair no segundo caso, o que torna o XML caro em \textit{tokens} em relação ao conteúdo que carrega. Como cada família de modelos tem seu próprio tokenizador, o mesmo texto pode ter números de \textit{tokens} diferentes em modelos diferentes.

O comprimento da sequência determina o custo de usar o modelo. Na atenção, cada posição se relaciona com todas as anteriores, e o custo cresce com o quadrado do comprimento. Os \textit{tokens} de entrada são processados em paralelo, de uma vez, mas os de saída são gerados um por passo, cada passo exigindo uma passagem pelo modelo inteiro. Por isso a saída é a parte mais lenta da inferência e, nos provedores que cobram por \textit{token}, a mais cara, com preços várias vezes maiores que os da entrada \cite{togetherpricing2026}. Alguns provedores oferecem também um preço menor para a parte da entrada que se repete entre requisições, como as instruções fixas de um \textit{prompt}, que fica guardada em \textit{cache} e não precisa ser processada de novo. O preço por \textit{token} cresce com o porte do modelo, e por isso a troca de um modelo de grande porte por um pequeno reduz o custo mesmo quando o número de \textit{tokens} é o mesmo \cite{belcak2025small}. A forma de cálculo do custo usada nos experimentos está na Seção~\ref{subsec:custo}.

A geração de artefatos estruturados por modelos de linguagem, como código-fonte, consultas e linguagens de marcação, tornou-se uma linha de pesquisa própria a partir de trabalhos como o de \citeonline{chen2021evaluating}, que avaliaram a geração de programas pela execução de testes. O problema central é que a geração probabilística não garante, por si só, que a saída siga uma sintaxe ou um esquema. Em um texto em linguagem natural, um desvio pequeno é tolerado pelo leitor. Em um arquivo XML, uma inconsistência local pode tornar todo o documento inválido.

\citeonline{poesia2022synchromesh} propõem combinar o modelo probabilístico com mecanismos determinísticos de restrição. Na decodificação restrita, a cada passo, os \textit{tokens} que levariam a uma saída fora da gramática são excluídos antes da escolha, e a saída fica garantidamente correta na sintaxe. Há outras duas formas de controle, que dispensam alterar a decodificação: validar a saída depois de gerada, rejeitando-a ou pedindo uma nova geração quando falha, e pedir ao modelo uma representação intermediária mais simples de gerar e de verificar, convertida depois no formato final por um programa. As três formas não são excludentes. A decodificação restrita garante a sintaxe, mas não reduz o número de \textit{tokens} gerados nem impede erros de conteúdo.

As referências cruzadas são um caso particular de dificuldade. Em formatos em que um elemento aponta para outro por um identificador, como os fluxos do XML do BPMN, a correção de uma referência depende de algo escrito em outro ponto do texto. Quando a referência aponta para trás, para um elemento já escrito, o identificador está no contexto do modelo e basta copiá-lo. Quando o formato exige marcar um elemento no momento em que ele é escrito, para que possa ser referenciado depois, o modelo precisa saber, naquele momento, que voltará a ele, o que exige a antecipação descrita anteriormente.

\section{Linguagens de Domínio Específico e Compressão Semântica}
\label{sec:dsl-compressao}\label{subsec:dsl}\label{subsec:gramaticas}\label{subsec:compressao-semantica}

Uma forma de reduzir o que o modelo precisa gerar é pedir uma notação mais curta que o formato final e convertê-la depois por programa. Esta seção apresenta as linguagens de domínio específico, que são o tipo de notação usado para isso, as gramáticas formais, que definem essas linguagens com precisão suficiente para que um programa as leia, e a ideia de compressão semântica, que orienta o que a notação curta deve conter.

Para \citeonline{fowler2010domain}, uma linguagem de domínio específico (DSL, do inglês \textit{Domain-Specific Language}) é uma linguagem de computador de expressividade limitada, voltada a um domínio particular. A limitação é intencional: por não precisar expressar qualquer programa, a linguagem pode usar os conceitos do próprio domínio e omitir detalhes que uma linguagem de propósito geral exigiria. \citeonline{mernik2005domain} apontam como benefícios recorrentes das DSLs a elevação do nível de abstração, a redução de código repetitivo e a facilidade de verificação e de análise. \citeonline{fowler2010domain} distingue ainda as DSLs internas, escritas dentro de uma linguagem de programação existente, das externas, que têm sintaxe própria e exigem um programa para lê-las.

A serialização XML da BPMN já é uma linguagem formal, mas foi concebida para a troca de arquivos entre ferramentas, e não para ser escrita por pessoas ou por modelos de linguagem. Grande parte do que ela contém, como os identificadores, os elementos de ligação e as declarações de espaços de nomes, pode ser derivada de uma descrição mais curta do processo por regras fixas. Uma DSL externa para processos pode, portanto, conter só o que não é derivável, e deixar o restante para um programa.

A sintaxe de uma linguagem de computador é definida por uma gramática formal. Nas gramáticas livres de contexto, usadas na maior parte das linguagens de programação, cada regra diz como um símbolo da linguagem pode ser decomposto em uma sequência de outros símbolos e de trechos literais de texto \cite{aho2006compilers}. A notação EBNF (\textit{Extended Backus-Naur Form}), padronizada pela ISO/IEC~14977 \cite{iso1996ebnf}, escreve essas regras de forma compacta, com operadores para repetição, opção e alternativa. A regra ilustrativa
\begin{equation*}
	\textit{fluxo} = \textit{elemento}, \; \{\, \texttt{\textquotesingle{}->\textquotesingle{}}, \; \textit{elemento} \,\} \, ;
\end{equation*}
diz que um fluxo é um elemento seguido de zero ou mais repetições do operador \texttt{->} e de outro elemento.

A partir da gramática, um analisador sintático lê um texto e produz a sua árvore sintática, que mostra como o texto se decompõe nas regras, ou rejeita o texto quando ele não segue a gramática. Esses analisadores podem ser gerados automaticamente a partir da própria gramática. O algoritmo de \citeonline{earley1970efficient} é um dos que aceitam qualquer gramática livre de contexto, inclusive as ambíguas, o que dá liberdade para escrever a gramática na forma mais legível. Sobre a árvore, um segundo programa produz a saída. Um programa que lê um texto em uma linguagem e o traduz em um texto equivalente em outra é um compilador, no sentido amplo de \citeonline{aho2006compilers}. Quando a saída é também um texto legível, e não código de máquina, usa-se o termo transpilador, ou tradutor de código-fonte para código-fonte.

Um transpilador escrito com regras fixas é determinístico: a mesma entrada produz sempre a mesma saída. Isso tem duas consequências para o uso com modelos de linguagem. A primeira é que a correção do transpilador pode ser verificada por programa, sobre todas as entradas disponíveis, independentemente de qualquer modelo. A segunda é que, se o transpilador está correto, toda saída do modelo aceita pela gramática produz um arquivo correto, e os erros que restam são os do próprio modelo, na sintaxe da DSL ou no conteúdo do processo.

Compressão semântica é a representação de um artefato por uma forma mais curta que contém apenas o que não pode ser derivado, de modo que a forma completa seja reconstruída por um programa determinístico. Ela difere da compressão de dados convencional, como a dos arquivos compactados, que opera sobre os \textit{bytes} sem considerar o significado e produz uma saída que não pode ser lida nem escrita por um modelo de linguagem. Difere também da compressão de \textit{prompts}, como a do LLMLingua \cite{jiang2023llmlingua}, que remove partes da entrada do modelo e aceita perder informação. Na compressão semântica, o que se reduz é a saída, e a redução não pode perder informação, porque a forma completa precisa ser reconstruída a partir dela.

No caso do BPMN, o que pode ser removido é o que o arquivo repete ou o que segue de regras: os identificadores, que podem ser gerados; os elementos de ligação, que podem ser dados pela ordem em que os nós aparecem; os \textit{gateways} de convergência, que podem ser inferidos do fim de um bloco de ramos; e as declarações de espaços de nomes, que são sempre as mesmas. O que precisa ficar é o conteúdo do processo: os nós, seus tipos e rótulos, a estrutura das decisões e das execuções paralelas e os pontos de retorno.

O ganho é medido pela razão entre o número de \textit{tokens} da forma completa e o da forma comprimida, contados com o mesmo tokenizador:
\begin{equation}
	\text{TCR} = \frac{\text{tokens}(\text{forma completa})}{\text{tokens}(\text{forma comprimida})}.
\end{equation}
Uma razão de 5, por exemplo, significa que a forma comprimida tem um quinto dos \textit{tokens}, ou seja, 80\% a menos. A forma como a razão é calculada nos experimentos está na Seção~\ref{subsec:tcr}. A razão, sozinha, não mede qualidade. Uma forma comprimida só tem valor se o modelo a gerar corretamente, e uma notação mais curta pode ser mais difícil de gerar do que o formato original, por exemplo se for mais distante dos textos vistos no treino do modelo. Por isso a compressão precisa ser avaliada junto com a validade das saídas.

\section{Ajuste de Modelos de Linguagem}
\label{sec:adaptacao}\label{subsec:sft}\label{subsec:lora}

Há duas formas de fazer um modelo de linguagem executar uma tarefa nova. Na primeira, a tarefa é descrita no \textit{prompt}, e o modelo a executa pelo aprendizado em contexto, sem alteração dos pesos (Seção~\ref{subsec:transformer}). Na segunda, os pesos são ajustados com exemplos da tarefa. A segunda forma permite usar um modelo menor e um \textit{prompt} mais curto, o que reduz o custo de cada requisição, e \citeonline{belcak2025small} argumentam que modelos pequenos ajustados são suficientes para tarefas repetitivas e bem delimitadas. Esta seção descreve o ajuste supervisionado e os métodos que o tornam viável com pouca memória.

No ajuste supervisionado (SFT, do inglês \textit{Supervised Fine-Tuning}), o treino de um modelo já treinado continua sobre pares de entrada e saída esperada. Para um par $(x, y)$, a perda é o negativo do logaritmo da probabilidade que o modelo atribui à saída esperada,
\begin{equation}
	\mathcal{L}(\theta) = -\sum_{t=1}^{T} \log p_{\theta}(y_t \mid y_{<t}, x),
\end{equation}
em que $\theta$ são os pesos do modelo. A soma percorre só os \textit{tokens} da saída, e os da entrada servem apenas de contexto. Minimizar essa perda aproxima a distribuição do modelo das saídas dos exemplos. \citeonline{ouyang2022training} usam esse procedimento como primeira etapa para fazer modelos seguirem instruções. Parte dos pares é separada para validação e não é usada no treino. A perda sobre esses pares indica se o modelo está aprendendo a tarefa ou apenas memorizando os exemplos de treino: quando a perda no treino continua caindo e a da validação passa a subir, o modelo começou a memorizar.

Quando as saídas dos pares são geradas por um modelo maior, o ajuste é uma forma de destilação. \citeonline{hinton2015distilling} propuseram treinar um modelo menor, o aluno, para reproduzir as distribuições de saída de um modelo maior, o professor. \citeonline{kim2016sequence} mostraram que, em modelos que geram sequências, basta treinar o aluno sobre as sequências produzidas pelo professor, sem acesso às suas distribuições, o que chamaram de destilação no nível da sequência. Esse é o caso quando os dados de treino são gerados por um modelo de grande porte acessado por um provedor. Nesse caso, o aluno herda também os erros do professor, e por isso os pares precisam ser filtrados antes do treino.

Ajustar todos os pesos de um modelo com bilhões de parâmetros exige guardar, para cada peso, o gradiente e o estado do otimizador, o que multiplica a memória necessária. O LoRA (\textit{Low-Rank Adaptation}), de \citeonline{hu2022lora}, congela os pesos originais e treina apenas uma correção de posto baixo para algumas matrizes do modelo. Para uma matriz de pesos $W_0 \in \mathbb{R}^{d \times k}$, a saída da camada passa a ser
\begin{equation}
	h = W_0 x + \frac{\alpha}{r} B A x,
\end{equation}
em que $B \in \mathbb{R}^{d \times r}$ e $A \in \mathbb{R}^{r \times k}$ são as matrizes treinadas, o posto $r$ é muito menor que $d$ e $k$, e $\alpha$ é um fator de escala. O número de parâmetros treinados por matriz cai de $d \cdot k$ para $r(d + k)$. A matriz $B$ começa com zeros, de modo que, no início do treino, o modelo ajustado é igual ao original. As matrizes treinadas formam um adaptador, que é guardado à parte e pode ser carregado sobre o modelo base no momento da inferência.

O QLoRA, de \citeonline{dettmers2023qlora}, reduz ainda mais a memória ao guardar os pesos congelados do modelo base em 4 \textit{bits}, em um formato próprio chamado \textit{NormalFloat}, e ao quantizar também as constantes da própria quantização, o que os autores chamam de dupla quantização. Os cálculos são feitos em 16 \textit{bits}, com os pesos convertidos no momento do uso, e só o adaptador é mantido e treinado em precisão maior. Os autores relatam desempenho próximo ao do ajuste completo em 16 \textit{bits}, com uma fração da memória.

\section{Avaliação de Modelos de Processo Gerados}
\label{sec:nl-bpmn}\label{sec:metricas}\label{subsec:precisao-recall}\label{subsec:juiz}

A maior parte do conhecimento sobre os processos de uma organização está em textos, como manuais e descrições de procedimentos, e não em modelos. A conversão automática desses textos em modelos de processo é estudada desde antes dos modelos de linguagem atuais. \citeonline{friedrich2011process} geravam modelos BPMN a partir de texto com técnicas tradicionais de Processamento de Linguagem Natural (PLN), como a análise sintática das frases, a resolução de correferências e regras que mapeavam as construções da língua em elementos do modelo. Os resultados mostraram que a tarefa é viável, mas também que os textos de processo deixam muita informação implícita.

Entre as dificuldades da tarefa, \citeonline{vanderaa2018challenges} apontam a ambiguidade das palavras e da estrutura das frases, a variação no nível de detalhe com que cada passo é descrito, as lacunas de informação, como ações sem responsável ou decisões sem a condição que as determina, e a existência de vários modelos válidos para uma mesma descrição. Este último ponto afeta diretamente a avaliação. Para muitos textos, não existe um único modelo correto, mas vários modelos compatíveis com o texto, que diferem na granularidade das atividades, na forma de representar uma decisão ou na interpretação de um trecho ambíguo.

Como consequência, um modelo de referência feito por uma pessoa é uma interpretação possível do texto, e não o único resultado correto. Um modelo gerado que difere da referência pode estar errado ou pode ser outra interpretação válida, e uma medida que compara o modelo gerado apenas com a referência não distingue os dois casos. Quando o texto é vago, quem modela precisa acrescentar detalhes que o texto não dá, e modeladores diferentes acrescentam detalhes diferentes. Os conjuntos de dados usados nos experimentos e os modelos de referência que eles oferecem estão descritos na Seção~\ref{sec:dados}.

A avaliação de um modelo de processo gerado responde a duas perguntas diferentes: se a saída é um arquivo correto e se ela representa o processo descrito no texto. A primeira pode ser respondida por programa. A segunda exige julgamento, feito por comparação com um modelo de referência ou por um avaliador, humano ou não. Além disso, quando várias configurações são avaliadas sobre os mesmos itens, é preciso decidir se as diferenças observadas entre elas são maiores que a variação do acaso, o que é tratado na Seção~\ref{sec:estatistica}.

A validade é verificada pela validação contra o esquema, descrita na Seção~\ref{subsec:bpmn-xml}. Ela é necessária, porque um arquivo rejeitado pelo esquema não pode ser aberto nem processado pelas ferramentas de modelagem. Não é suficiente, porque o esquema não verifica o comportamento do processo, e um arquivo válido pode descrever o processo errado. A fidelidade ao texto precisa, assim, ser medida à parte.

A forma mais comum de medir a fidelidade é comparar o modelo gerado com um modelo de referência. Os dois modelos são reduzidos a conjuntos de itens comparáveis, como os rótulos das atividades ou os pares de nós ligados por um fluxo, e a concordância é medida pela precisão, pela revocação e pela medida F1. Sendo $G$ o conjunto de itens do modelo gerado e $E$ o do modelo de referência,
\begin{equation}
	P = \frac{|G \cap E|}{|G|}, \qquad
	R = \frac{|G \cap E|}{|E|}, \qquad
	F_1 = \frac{2 \, P \, R}{P + R}.
\end{equation}
A precisão penaliza o que o modelo gerado tem a mais, a revocação penaliza o que ele tem a menos, e a F1 é a média harmônica das duas. Como os rótulos raramente coincidem letra a letra, a comparação exige antes alinhar os nós dos dois modelos, por igualdade ou por semelhança dos rótulos.

Medidas mais elaboradas comparam a estrutura ou o comportamento dos modelos. \citeonline{dijkman2011similarity} comparam modelos de processo pela semelhança dos rótulos, pela estrutura do grafo e pelo comportamento que os modelos permitem. A comparação estrutural usa a distância de edição em grafos (GED, do inglês \textit{Graph Edit Distance}), definida como o custo mínimo de uma sequência de inserções, remoções e substituições de nós e de arestas que transforma um grafo no outro \cite{sanfeliu1983distance}. O cálculo exato da GED é computacionalmente difícil, e na prática usam-se aproximações \cite{gao2010graph}. \citeonline{weidlich2011behavioral} propõem comparar os modelos pelas relações de ordem entre pares de atividades, como precedência estrita, exclusão e intercalação, o que torna a medida indiferente a formas diferentes de desenhar o mesmo comportamento.

Todas essas medidas dependem do modelo de referência e herdam o problema descrito no início desta seção: um modelo gerado que represente outra interpretação válida do texto é penalizado como se estivesse errado. Dependem também do alinhamento dos rótulos, e uma paráfrase como ``Verify invoice'' no lugar de ``Check invoice'' pode ser contada como uma atividade a mais e outra a menos.

Uma alternativa à comparação com a referência é pedir a um modelo de linguagem que julgue a saída diretamente contra o texto de origem. \citeonline{zheng2023judging} avaliaram essa prática, que chamaram de LLM como juiz (\textit{LLM-as-a-judge}), e encontraram concordância entre os juízes mais capazes e avaliadores humanos semelhante à concordância entre dois humanos, na avaliação de respostas a perguntas abertas. Os mesmos autores descrevem vieses desses juízes: a preferência pela resposta apresentada em certa posição, a preferência por respostas mais longas e a preferência por respostas do próprio modelo. \citeonline{panickssery2024llm} mostraram que essa última está ligada à capacidade do juiz de reconhecer o próprio texto: quanto melhor o modelo distingue as próprias saídas, mais ele as favorece.

Esses vieses orientam o desenho de um juiz. Usar um juiz de família diferente da dos modelos avaliados reduz a preferência pelas próprias saídas. Apresentar todas as saídas na mesma forma, sem indicar de onde vieram, impede que o juiz reconheça a origem pelo formato. Uma rubrica com critérios explícitos e a exigência de que o juiz liste os problemas encontrados antes de dar a nota tornam cada julgamento auditável. Repetir o julgamento e resumir as notas pela mediana reduz o efeito da variação entre execuções.

A concordância observada em uma tarefa não garante concordância em outra, e um juiz precisa ser validado na tarefa em que será usado. A forma mais direta é comparar as suas notas com as de especialistas sobre os mesmos itens. Como as notas são ordinais, a comparação usa a correlação de postos de \citeonline{spearman1904proof}, que mede se os dois avaliadores ordenam os itens da mesma forma, sem exigir que usem a escala do mesmo jeito. Sem empates, ela é
\begin{equation}
	\rho = 1 - \frac{6 \sum_{i=1}^{n} d_i^2}{n (n^2 - 1)},
\end{equation}
em que $d_i$ é a diferença entre os postos do item $i$ nas duas avaliações e $n$ é o número de itens. A segunda forma é testar o comportamento do juiz com casos construídos, à maneira dos testes de comportamento de \citeonline{ribeiro2020beyond}. Em um teste de invariância, a saída é alterada de um modo que não muda o seu significado, como a troca de rótulos por sinônimos, e a nota não deve mudar. Em um teste direcional, introduz-se um defeito conhecido, como a remoção de uma atividade, e a nota deve cair. Os dois testes não dependem de notas humanas, e por isso podem ser aplicados mesmo quando as notas disponíveis são poucas ou pouco confiáveis.

\section{Comparação Pareada e Testes de Equivalência}
\label{sec:estatistica}\label{subsec:estatistica}

Quando todas as configurações são avaliadas sobre os mesmos itens, a comparação entre duas delas pode ser pareada. Para cada item $i$, calcula-se a diferença $d_i$ entre os valores das duas configurações, e a comparação é feita sobre a média dessas diferenças. O pareamento remove da comparação a variação de dificuldade entre os itens, que é comum às duas configurações, e por isso torna a comparação mais precisa do que comparar as médias de cada uma.

O \textit{bootstrap}, proposto por \citeonline{efron1979bootstrap}, estima a incerteza dessa média sem supor uma distribuição para as diferenças. Sorteiam-se, com reposição, $n$ itens entre os $n$ disponíveis, e calcula-se a média das diferenças da amostra sorteada. Repetido $B$ vezes, o procedimento produz $B$ médias, e o intervalo de confiança de 95\% pelo método dos percentis vai do percentil 2,5 ao percentil 97,5 dessas médias \cite{efron1993introduction}. O teste de postos sinalizados de \citeonline{wilcoxon1945individual} verifica se as diferenças pendem para um dos lados mais do que o acaso explicaria. Ele ordena os valores absolutos das diferenças, soma os postos das diferenças positivas e compara essa soma com a esperada se as configurações fossem iguais, também sem supor uma distribuição normal. Quando vários contrastes são testados, a chance de pelo menos um resultado falso cresce com o número de testes, e a correção de \citeonline{holm1979simple} ajusta os níveis de significância para controlar essa chance.

Um teste sem diferença significativa não mostra que as configurações são iguais: com poucos itens, o intervalo de confiança é largo e pode conter tanto o zero quanto diferenças grandes \cite{altman1995absence}. Para afirmar igualdade, é preciso um teste de equivalência. No procedimento de dois testes unilaterais (TOST, do inglês \textit{Two One-Sided Tests}), de \citeonline{schuirmann1987comparison}, fixa-se antes uma margem $\Delta$, a maior diferença considerada irrelevante na prática, e as configurações são declaradas equivalentes quando o intervalo de confiança da diferença está inteiro dentro de $[-\Delta, \Delta]$ \cite{lakens2017equivalence}. Combinado ao teste usual, o intervalo leva a três conclusões possíveis: equivalência, quando está inteiro dentro da margem; diferença além da margem, quando está inteiro fora; e resultado inconclusivo, nos demais casos. No TOST com nível de 5\%, usa-se o intervalo de 90\%, e exigir que o intervalo de 95\% esteja dentro da margem é, portanto, um critério mais rigoroso. A margem precisa ser justificada sem recorrer aos dados comparados, porque uma margem escolhida depois de vistos os resultados pode ser ajustada para produzir a conclusão desejada.
